%% =====================================================================
%%  La collection : les cinq tomes de la méthode
%% =====================================================================
\liminaire{collection}{La collection}{De zéro à l'aisance : cinq tomes, cinq niveaux du CECRL}

\newcommand{\otCarteTome}[8]{% n, niveau, couleur, titre, sous-titre, savoir-faire, grammaire, statut
  \begin{tcolorbox}[enhanced,colback=white,colframe=otRule,boxrule=0.5pt,arc=2mm,left=0mm,right=4mm,
    top=0mm,bottom=0mm,boxsep=0pt,before skip=0pt,after skip=3.2mm,
    overlay={\fill[#3,rounded corners=2mm] (frame.north west) rectangle ([xshift=30mm]frame.south west);
      \fill[#3] ([xshift=26mm]frame.north west) rectangle ([xshift=30mm]frame.south west);
      \node[text=white,font=\otSemi\fontsize{7}{8}\selectfont,anchor=north] at ([xshift=15mm,yshift=-4.5mm]frame.north west) {\otCaps[12]{Tome}};
      \node[text=white,font=\otBold\fontsize{30}{30}\selectfont] at ([xshift=15mm,yshift=-15.5mm]frame.north west) {#1};
      \node[fill=white,text=#3,rounded corners=1.2mm,inner xsep=2mm,inner ysep=1mm,font=\otBold\fontsize{10}{11}\selectfont]
        at ([xshift=15mm,yshift=-26mm]frame.north west) {#2};}]
  \hspace*{34mm}\begin{minipage}[t]{\dimexpr\linewidth-34mm\relax}\vspace{3.6mm}
    {\otSemi\fontsize{14}{16}\selectfont #4}\hspace{2.5mm}{\otRegular\fontsize{9.4}{12}\selectfont\color{otMuted}#5}\hfill
    {\otSemi\fontsize{7}{8.5}\selectfont\color{#3}\otCaps[10]{#8}}\par\vspace{1.6mm}
    {\fontsize{9.8}{12.4}\selectfont #6\par}\vspace{1.2mm}
    {\fontsize{9}{11.4}\selectfont\color{otMuted}{\otSemi\fontsize{7}{8.5}\selectfont\color{otViolet}\otCaps[10]{Grammaire}}\enspace #7\par}
    \vspace{3.6mm}
  \end{minipage}
  \end{tcolorbox}}

\otCarteTome{1}{A1}{otGreen}{Premiers pas}{Débutant}
  {Se présenter, parler de son quotidien, se débrouiller en ville, raconter un week-end. Avec Léa,
   à New York, pendant ses premiers mois.}
  {\en{be} · présent simple · présent en \en{-ing} · \en{can} · prétérit · \en{be going to}}{Ce livre}
\otCarteTome{2}{A2}{otGreenDark}{Les bases solides}{Élémentaire}
  {Échanger sur des sujets familiers, comparer, raconter ses expériences, parler de l'avenir,
   voyager en autonomie. Léa découvre Londres avec Emma.}
  {\en{present perfect} · comparatifs et superlatifs · \en{will} · \en{past continuous} · \en{must}, \en{have to}, \en{should}}{À paraître}
\otCarteTome{3}{B1}{otNavy!75}{Communiquer avec aisance}{Intermédiaire}
  {Donner son avis et le justifier, raconter une histoire, faire face aux imprévus, travailler en
   anglais au quotidien.}
  {\en{present perfect continuous} · \en{past perfect} · conditionnels · voix passive · discours indirect · relatives}{À paraître}
\otCarteTome{4}{B2}{otNavy}{Parler couramment}{Avancé}
  {Argumenter, comprendre films, séries et presse, mener une réunion, rédiger des e-mails
   professionnels.}
  {conditionnels passés · \en{wish}, \en{if only} · modaux au passé · \en{phrasal verbs} · mots de liaison}{À paraître}
\otCarteTome{5}{C1}{otNavyDark}{Maîtriser l'anglais}{Autonome}
  {Nuancer sa pensée, adapter son registre, saisir l'implicite et l'humour, s'exprimer avec
   précision dans tous les contextes.}
  {structures d'insistance · subjonctif · expressions idiomatiques · collocations · écrit formel}{À paraître}

\begin{astuce}[Quel tome choisir ?]
Vous débutez, ou vos souvenirs d'anglais sont très lointains : commencez par ce tome 1.
Vous comprenez les phrases simples mais butez dès que la conversation s'accélère : le tome 2.
En cas de doute, faites les tests express des unités du tome 1 : si vous obtenez plus de 8 sur 10
partout, passez au tome suivant.
\end{astuce}
